\documentclass[12pt,a4paper]{exam}

% --------------------
% Packages
% --------------------
\usepackage[utf8]{inputenc}
\usepackage[T1]{fontenc}
\usepackage[french]{babel}
\usepackage{amsmath,amssymb}
\usepackage{geometry}
\usepackage{pgfplots}
\pgfplotsset{compat=1.18}
\geometry{margin=2cm}

% --------------------
% En-tête / pied (exam)
% --------------------
\header{}{DS de mathématiques}{1ère STMG}
\footer{}{\thepage}{}

% --------------------
\begin{document}

\begin{center}
    \Large\textbf{Interrogation\\Dérivation}
\end{center}

\vspace{0.3cm}

\noindent
\textbf{Durée :} 1h \hfill
\textbf{Calculatrice :} \underline{NON} autorisée

\vspace{0.5cm}

% ======================================================
\noindent
\textbf{Exercice 1 : Dérivées et tangentes (\ \ \ \  / 5 points )}
\vspace{10pt}
\noindent\\On considère la fonction $f$ représentée ci-dessous ainsi que deux tangentes à $f$.

\vspace{-5pt}
\begin{center}
    \begin{tikzpicture}[scale=1, xscale=1.5]

    % Grille
    \draw[step=1cm, gray, very thin] (-3,-2) grid (6,7);

    % Axes
    \draw[->] (-3.2,0) -- (6.2,0) node[right] {$x$};
    \draw[->] (0,-2.2) -- (0,7.2) node[above] {$y$};

    % Graduations x
    \foreach \x in {-3,-2,-1,1,2,3,4,5,6}
        \draw (\x,0.1) -- (\x,-0.1) node[below] {\small $\x$};

    % Graduations y
    \foreach \y in {-1,1,2,3,4,5,6}
        \draw (0.1,\y) -- (-0.1,\y) node[left] {\small $\y$};

    % Courbe y = 3x^2 - 4 (limitée à y ≤ 5)
    \draw[domain=-1.5:3.5, samples=200, smooth, thick, black]
        plot (\x, {\x*\x - 2*\x + 2});

    \node at (-1,6.2) {$\mathcal{C}_f$};

    % Courbe y = 3x^2 - 4 (limitée à y ≤ 5)
    \draw[domain=-2.5:2, samples=200, smooth, very thick, black]
        plot (\x, {-2*\x + 2});

    \node at (-1.5,4.5) {$\mathcal{T}_A$};

    % Courbe y = 3x^2 - 4 (limitée à y ≤ 5)
    \draw[domain=1.25:3.5, samples=200, smooth, very thick, black]
        plot (\x, {4*\x -7});

    \node at (2.5,1.75) {$\mathcal{T}_B$};

    \filldraw (0, 2) circle (2pt) node[above right] {A};

    \filldraw (3, 5) circle (2pt) node[below right] {B};

    \end{tikzpicture}
\end{center}


\begin{questions}

\question[2] Déterminer graphiquement $f'(0)$ et $f'(3)$.

\question[3] En déduire les équations des tangentes $\mathcal{T}_A$ et $\mathcal{T}_B$ à 
la courbe $\mathcal{C}_f$.

\end{questions}

\vspace{1cm}

% ========================================================
%             Exercice 2
% ========================================================
\noindent
\textbf{Exercice 2 : Calculs de dérivées ( \ \ \ / 5 points )}

\vspace{10pt}
\noindent

\begin{questions}

\question[2] Déterminer l'expression de la fonction dérivée de chacune des fonctions suivantes :

\vspace{5pt}
\hspace{1cm}a) $f(x) = 3x^3 - 7x$ \hspace{1cm}b) $g(x) = 3x^2 - x + 5$

\vspace{10pt}

\question On considère la fonction : $f(x) = 2x^2 + 4x - 2$.

\vspace{5pt}
\hspace{1cm}a) (1 point) Calculer la dérivée $f'(x)$.

\vspace{5pt}
\hspace{1cm}b) (1 point) Étudier le signe de cette dérivée.

\vspace{5pt}
\hspace{1cm}c) (1 point) Dresser le tableau de variation de la fonction $f$.

\end{questions}

\vspace{0.5cm}

\begin{center}
    \textit{Fin du devoir}
\end{center}

\end{document}
